\documentclass[11pt,a4paper]{article}
\usepackage[utf8]{inputenc}
\usepackage[T1]{fontenc}
\usepackage{amsmath,amssymb}
\usepackage[french]{babel}
\frenchbsetup{StandardLists=true}
\usepackage{enumitem}
\usepackage[left=2cm,right=2cm,top=2.5cm,bottom=2cm]{geometry}
\usepackage{multirow,tabularx,booktabs}
\usepackage[dvipsnames]{xcolor}
\usepackage{fouriernc}
\usepackage{fancyhdr}
\usepackage{colortbl}
\usepackage{pifont}
\usepackage[np]{numprint}

\pagestyle{fancy}
\renewcommand\headrulewidth{1pt}
\fancyhead[L]{Lycée Denis Diderot}
\fancyhead[R]{1BTS}
\fancyfoot[C]{page \thepage}
\fancyfoot[R]{Année 2025-2026}
\fancyfoot[L]{Alexis RUIZ}

%------ Couleurs ------
\definecolor{exoheader}{RGB}{30,90,160}
\definecolor{exolight}{RGB}{220,235,255}
\definecolor{solcolor}{RGB}{180,0,0}
\newcommand{\sol}[1]{\textcolor{solcolor}{\textbf{#1}}}
\newcommand{\solc}[1]{\textcolor{solcolor}{#1}}

%------ Titre d'exercice sans tcolorbox ------
\newcommand{\titreexo}[2]{%
  \vspace{4mm}%
  \noindent\colorbox{exoheader}{\color{white}\bfseries\strut\enspace #1\enspace}%
  \hfill\textbf{#2}%
  \vspace{1mm}\par\noindent\textcolor{exoheader}{\rule{\linewidth}{1pt}}\par
  \vspace{2mm}%
}

\renewcommand{\arraystretch}{1.6}
\setlength{\parindent}{0pt}
\newcolumntype{L}[1]{>{\raggedright\arraybackslash}p{#1}}

% ============================================================
\begin{document}

\begin{center}
  \Large\textcolor{exoheader}{\textbf{GRILLE DE CORRECTION}}\\[2mm]
  \large Évaluation \textit{-- Probabilités conditionnelles --} 1BTS\\[4mm]
 
\end{center}

\vspace{2mm}\hrule\vspace{3mm}

% ============================================================
\titreexo{Exercice 1 \,--\, Vocabulaire}{/ 4 points}

\textbf{Partie A -- Vrai ou Faux}
\hfill\small\textit{0,5 pt par bonne réponse \,--\, sous-total : \textbf{2 pts}}

\medskip
\begin{tabularx}{\linewidth}{|>{\raggedright\arraybackslash}X|>{\centering\arraybackslash}p{1.5cm}|>{\raggedright\arraybackslash}p{5cm}|>{\centering\arraybackslash}p{1.5cm}|}
\hline
\rowcolor{exolight}
\textbf{Affirmation} & \textbf{V / F} & \textbf{Correction si fausse} & \textbf{Pts}\\
\hline
Pour tout événement $A$ : $P(\bar{A}) = 1 + P(A)$
  & \sol{F} & \solc{$P(\bar{A}) = 1 - P(A)$} & 0,5\\
\hline
Si $A$ et $B$ sont incompatibles alors $P(A \cup B) = P(A) + P(B)$
  & \sol{V} & & 0,5\\
\hline
$P_B(A) = \dfrac{P(A \cap B)}{P(B)}$
  & \sol{V} & & 0,5\\
\hline
Si $A$ et $B$ sont indépendants alors $P(A \cap B) = P(A) + P(B)$
  & \sol{F} & \solc{$P(A\cap B) = P(A)\times P(B)$} & 0,5\\
\hline
\multicolumn{3}{|r|}{\textbf{Sous-total Partie A}} & \sol{\textbf{2,0}}\\
\hline
\end{tabularx}

\smallskip

\textbf{Partie B -- Associer}
\hfill\small\textit{0,5 pt par bonne association \,--\, sous-total : \textbf{2 pts}}

\medskip
\begin{tabularx}{\linewidth}{|>{\centering\arraybackslash}p{2.8cm}|>{\raggedright\arraybackslash}X|>{\centering\arraybackslash}p{1.5cm}|}
\hline
\rowcolor{exolight}
\textbf{Notation} & \textbf{Association correcte} & \textbf{Pts}\\
\hline
$P(A \cup B)$ & \solc{Probabilité que $A$ \textbf{ou} $B$ se réalise} & 0,5\\
\hline
$P(A \cap B)$ & \solc{Probabilité que $A$ \textbf{et} $B$ se réalisent simultanément} & 0,5\\
\hline
$P_B(A)$      & \solc{Probabilité de $A$ \textbf{sachant} $B$} & 0,5\\
\hline
$P(A)$        & \solc{Probabilité que $A$ se réalise} & 0,5\\
\hline
\multicolumn{2}{|r|}{\textbf{Sous-total Partie B}} & \sol{\textbf{2,0}}\\
\hline
\end{tabularx}

% ============================================================
\titreexo{Exercice 2 \,--\, Arbre à compléter}{/ 4 points}

\textit{Urne : 4 boules rouges, 6 bleues -- tirage sans remise.}

\medskip
\textbf{Q1. Branches manquantes} \hfill\small\textit{0,5 pt par valeur correcte \,--\, sous-total : \textbf{2 pts}}

\medskip
\begin{tabularx}{\linewidth}{|>{\raggedright\arraybackslash}X|>{\centering\arraybackslash}p{4.5cm}|>{\centering\arraybackslash}p{1.5cm}|}
\hline
\rowcolor{exolight}
\textbf{Branche / probabilité} & \textbf{Valeur attendue} & \textbf{Pts}\\
\hline
$P(B)$ (branche bleue, 1\textsuperscript{er} tirage) & \sol{$\dfrac{6}{10} = \dfrac{3}{5}$} & 0,5\\
\hline
$P_{R}(R)$ (rouge $|$ rouge au 1\textsuperscript{er} tirage) & \sol{$\dfrac{3}{9} = \dfrac{1}{3}$} & 0,5\\
\hline
$P_{B}(B)$ (bleue $|$ bleue au 1\textsuperscript{er} tirage) & \sol{$\dfrac{5}{9}$} & 0,5\\
\hline
Probabilités des 4 feuilles & \solc{$P(R,R)=\frac{2}{15}$ \quad $P(R,B)=\frac{4}{15}$ \quad $P(B,R)=\frac{4}{15}$ \quad $P(B,B)=\frac{1}{3}$} & 0,5\\
\hline
\multicolumn{2}{|r|}{\textbf{Sous-total Q1}} & \sol{\textbf{2,0}}\\
\hline
\end{tabularx}

\smallskip

\textbf{Q2.} Probabilité d'obtenir deux boules rouges \hfill\textit{(1 pt)}

\medskip
\sol{$P(R,R) = \dfrac{2}{5} \times \dfrac{1}{3} = \dfrac{2}{15}$}
\hfill\textit{Calcul correct : 0,5 pt \,+\, résultat : 0,5 pt}

\smallskip

\textbf{Q3.} Probabilité d'obtenir au moins une boule rouge \hfill\textit{(1 pt)}

\medskip
\sol{$P(\text{au moins un rouge}) = 1 - P(B,B) = 1 - \dfrac{1}{3} = \dfrac{2}{3}$}
\hfill\textit{Méthode complémentaire ou somme directe acceptée : 0,5 + 0,5}

% ============================================================
\titreexo{Exercice 3 \,--\, Construction et utilisation d'un arbre}{/ 6 points}

\textit{Fournisseurs $S_1$ (40\,\%) et $S_2$ (60\,\%) -- taux de défectueux : 4\,\% et 2\,\%.}

\medskip
\begin{tabularx}{\linewidth}{|>{\centering\arraybackslash}p{5mm}|>{\raggedright\arraybackslash}X|>{\raggedright\arraybackslash}p{6cm}|>{\centering\arraybackslash}p{1.5cm}|}
\hline
\rowcolor{exolight}
\textbf{Q} & \textbf{Éléments attendus} & \textbf{Résultat / commentaire} & \textbf{Pts}\\
\hline
1
  & Arbre correctement construit avec toutes les branches et probabilités étiquetées
  & \solc{$P(S_1)=0{,}4$, $P(S_2)=0{,}6$, $P_{{S_1}}(D)=0{,}04$, $P_{{S_1}}(\bar{D})=0{,}96$, $P_{{S_2}}(D)=0{,}02$, $P_{{S_2}}(\bar{D})=0{,}98$}
  & 2\\
\hline
2
  & $P(S_1\cap D) = P(S_1)\times P_{S_1}(D)$ et idem pour $S_2$
  & \sol{$P(S_1\cap D)=0{,}016$} \newline \sol{$P(S_2\cap D)=0{,}012$}
  & 1\\
\hline
3
  & Formule des probabilités totales : $P(D)=P(S_1\cap D)+P(S_2\cap D)$
  & \sol{$P(D) = 0{,}016 + 0{,}012 = 0{,}028$}
  & 1\\
\hline
\multirow{2}{*}{4}
  & Formule de Bayes : $P_D(S_1) = \dfrac{P(S_1\cap D)}{P(D)}$
  & \sol{$P_D(S_1) = \dfrac{0{,}016}{0{,}028} = \dfrac{4}{7} \approx 0{,}571$}
  & 1\\
\cline{2-4}
  & Interprétation correcte en français
  & \solc{Parmi les composants défectueux, environ 57,1\,\% proviennent de $S_1$}
  & 1\\
\hline
\multicolumn{3}{|r|}{\textbf{Total Exercice 3}} & \sol{\textbf{6}}\\
\hline
\end{tabularx}

% ============================================================
\titreexo{Exercice 4 \,--\, Tableau de probabilités}{/ 6 points}

\textit{$\np{2000}$ patients -- Service $A$ : 60\,\% avec 8\,\% de complications -- Service $B$ : 40\,\% avec 5\,\% de complications.}

\medskip
\textbf{Q1. Tableau complété} \hfill\small\textit{0,5 pt par case correcte \,--\, sous-total : \textbf{2 pts}}

\medskip
\begin{tabularx}{\linewidth}{|>{\raggedright\arraybackslash}X|>{\centering\arraybackslash}p{2.8cm}|>{\centering\arraybackslash}p{2.8cm}|>{\centering\arraybackslash}p{2.8cm}|}
\hline
\rowcolor{exolight}
 & \textbf{Service $A$} & \textbf{Service $B$} & \textbf{Total}\\
\hline
Avec complication $C$       & \sol{96}    & \sol{40}   & \sol{136}\\
\hline
Sans complication $\bar{C}$ & \sol{1\,104} & \sol{760}  & \sol{1\,864}\\
\hline
Total                        & \sol{1\,200} & \sol{800}  & $\np{2000}$\\
\hline
\end{tabularx}

\smallskip
\solc{\small $A$ : $60\,\%\times\np{2000}=\np{1200}$ \quad avec $C$ : $8\,\%\times\np{1200}=96$ \qquad
      $B$ : $40\,\%\times\np{2000}=800$ \quad avec $C$ : $5\,\%\times 800=40$}

\medskip
\begin{tabularx}{\linewidth}{|>{\centering\arraybackslash}p{5mm}|>{\raggedright\arraybackslash}X|>{\raggedright\arraybackslash}p{5.5cm}|>{\centering\arraybackslash}p{1.5cm}|}
\hline
\rowcolor{exolight}
\textbf{Q} & \textbf{Éléments attendus} & \textbf{Résultat} & \textbf{Pts}\\
\hline
2
  & $P(C) = \dfrac{\text{total }C}{\text{total}}$
  & \sol{$P(C) = \dfrac{136}{\np{2000}} = 0{,}068$}
  & 1\\
\hline
\multirow{2}{*}{3}
  & Définition de $A \cap C$ en français
  & \solc{Patient dans service $A$ \textbf{et} présentant une complication}
  & 0,5\\
\cline{2-4}
  & Calcul de $P(A \cap C)$
  & \sol{$P(A\cap C) = \dfrac{96}{\np{2000}} = 0{,}048$}
  & 0,5\\
\hline
\multirow{2}{*}{4}
  & $P_C(A) = \dfrac{P(A\cap C)}{P(C)}$
  & \sol{$P_C(A) = \dfrac{0{,}048}{0{,}068} = \dfrac{12}{17} \approx 0{,}706$}
  & 1\\
\cline{2-4}
  & Interprétation correcte en français
  & \solc{Parmi les patients avec complication, environ 70,6\,\% sont en service $A$}
  & 1\\
\hline
\multicolumn{3}{|r|}{\textbf{Total Exercice 4}} & \sol{\textbf{6}}\\
\hline
\end{tabularx}

\vspace{6mm}

\begin{center}
  \colorbox{exolight}{%
    \parbox{8cm}{\centering
      \vspace{3mm}
      \textbf{\textcolor{exoheader}{Note finale : \rule{2.5cm}{0.4pt} / 20 points}}\\[3mm]
    }%
  }
\end{center}

\end{document}